\section{Conclusion}
\label{sec:conclusion}

In this work, we presented \textbf{C-Lie-MM}, a mathematically rigorous, differentiable Riemannian-geometric framework for dynamic model ensembling and representation-space projection operator merging. We addressed a fundamental limitation of existing flat merging and routing methods: when blending low-rank projection operators, linear interpolation violates the geometry of the projection manifold, causing projected coordinate collapse and severe representation norm decay due to eigenvalue shrinkage.

By treating task subspaces as points on the Grassmannian manifold $\mathcal{G}(d, D)$, C-Lie-MM leverages formal Riemannian-geometric operations—specifically, the Grassmannian logarithm and exponential maps relative to a fixed, pre-computed reference point $Y_0$—to perform weighted ensembling along geodesics in the tangent space. This guarantees that the merged projection operator is always symmetric, idempotent, and of rank $d$. Furthermore, our fixed-reference formulation completely resolves the step-function discontinuity of dynamic barycenters (argmax jumps), restoring full differentiability, while enabling offline pre-computation of logarithms to deliver a massive $K$-times speedup in forward pass SVD operations.

Our experiments in a $14$-layer Analytical Coordinate Sandbox show that C-Lie-MM completely resolves the representation collapse of uniform merging under severe task interference, achieving $70.30\% \pm 4.01\%$ accuracy, and remains perfectly immune to heterogeneity collapse on streaming workloads. We believe that our geometric perspective on subspace ensembling opens a promising avenue for multi-task deep learning. A particularly exciting direction for future work is scaling C-Lie-MM to massive Generative Pre-trained Transformers (such as LLaMA-3 or Mistral-7B) and Vision Transformers (ViTs) for multi-task LoRA merging. Evaluating these real-world models on standard benchmarks like GLUE and GSM8K will allow us to compile actual physical weight evaluation numbers on physical weight tensors, comparing the downstream accuracy of C-Lie-MM against state-of-the-art parameter merging baselines such as TIES-Merging and ZipIt. This physical evaluation will directly complement our current high-fidelity simulations by validating the framework's effectiveness in real-world settings. To achieve this, we are designing a concrete, actionable roadmap to integrate C-Lie-MM into physical pre-trained weights:
\begin{enumerate}
    \item \textbf{Hugging Face PEFT Integration:} We plan to implement C-Lie-MM as a custom merge-method wrapper within the Hugging Face PEFT library. The task-specific projection bases $V_k$ will be extracted dynamically by performing an offline, localized SVD on the down-projection weight matrices (e.g., the $W_A$ matrices of the query and value LoRA layers).
    \item \textbf{Token-Level and Sequence-Level Serving Hook:} For encoder-only models (such as RoBERTa-Large) on GLUE, we hook the Gibbs routing policy to evaluate task routing coefficients sequence-wise using the mean-pooled token hidden activations at intermediate transformer layers. For autoregressive decoder-only models (such as LLaMA-3) on GSM8K, prompt-level mean-pooled activations are used to compute and lock the routing coefficients, ensuring zero decoding latency overhead.
    \item \textbf{End-to-End Hardware Profiling:} We plan to conduct exhaustive end-to-end latency and throughput benchmarks using our custom fused Triton GPU kernel on NVIDIA H100 and A100 GPU clusters, comparing physical serving performance against standard LoRA merging baselines (such as TIES-Merging and ZipIt). Specifically, we will profile our custom kernel to measure absolute latency speedups, SRAM-level caching efficiency, and token-generation throughput under varying batch sizes ($B$) and token context lengths, establishing a direct serving baseline against standard Hugging Face PEFT multi-LoRA serving.
\end{enumerate}
This real-world vision/NLP integration roadmap will bridge the gap between our high-fidelity sequential simulators and massive industrial transformer workloads. Additionally, future research will explore extending this framework to other matrix manifolds, such as the Stiefel manifold $\mathcal{V}_k(\mathbb{R}^n)$ for partial orthonormal structures. Finally, because our routing coefficients $\alpha_k$ are continuously differentiable, we can leverage them during joint training as a real-time signal for active task-difficulty discovery and automatic curriculum selection. By monitoring the routing distribution across tasks, we can dynamically adjust the sampling frequency or loss weighting of different tasks, potentially accelerating multi-task optimization and sample efficiency.
